% !PRINTED 231
\weberpage{231}
%<<CONT
and find by a simple calculation:
\weq{3}{H = -(x_1^{20} + x_2^{20}) + 228\,(x_1^{15}x_2^5 - x_2^{15}x_1^5) - 494\,x_1^{10}x_2^{10}.}
The fundamental form of degree 30 we can define as the functional determinant of $H$ and $f$. If we set
\weq{4}{T = \frac{1}{20}\,[f'(x_1)\,H'(x_2) - f'(x_2)\,H'(x_1)],}
then one finds
\weq{5}{T = (x_1^{30} + x_2^{30}) + 522\,(x_1^{25}x_2^5 - x_1^5x_2^{25}) - 10005\,(x_1^{20}x_2^{10} + x_2^{20}x_1^{10}).}
There is also an identity among these three forms. To find it, we set
\[
\lambda = x_1^{10} - x_2^{10}, \quad \mu = x_1^5 x_2^5,
\]
and obtain, if we take out the factor $x_1^{10} + x_2^{10}$ from $T$ and then form $T^2$:
\[
f^5 = \mu\,(\lambda + 11\mu)^5,
\]
\[
H = -\lambda^2 + 228\,\lambda\mu - 496\,\mu^2,
\]
\[
T^2 = (\lambda^2 + 4\mu^2)\,(\lambda^2 + 522\,\lambda\mu - 10004\,\mu^2)^2.
\]
From this one obtains by numerical computation
\weq{6}{T^2 + H^3 = 1728\,f^5,}
which is the relation sought.

Since the icosahedral group, as we have seen, is simple, it can by \S.\,40 have no relative invariants. From this it follows that the icosahedral forms $f$, $H$, $T$ remain absolutely unchanged when one subjects them to the binary icosahedral substitutions represented with determinant $1$. But the same can also be demonstrated directly in the following way, without making use of those general theorems.

The relation (6) can be presented in the form:
\weq{7}{\frac{T^2}{f^5} + \frac{H^3}{f^5} = 1728,}
and the two quotients $T^2 : f^5$, $H^3 : f^5$ we can regard as functions of the single variable $x = x_1 : x_2$. If we apply to this variable a substitution of the icosahedral group, these two quotients change only by constant factors; that is, if we
\weq{8}{\frac{T^2}{f^5} = \Phi(x), \quad \frac{H^3}{f^5} = \Psi(x), \quad y = \frac{\alpha x + \beta}{\gamma x + \delta}}
%CONT>>
